% Content of the Rudin Functional Analysis blueprint.
% Each \chapter mirrors one chapter of Rudin, Functional Analysis (2nd ed.).

\chapter{Topological Vector Spaces}

\begin{theorem}[Rudin 1.6]
  \label{thm:translation-homeo}
  \lean{Rudin.Ch01.translationHomeo, Rudin.Ch01.smulHomeo}
  \leanok
  Let $X$ be a topological vector space over $\mathbb{K}$.
  For each $a \in X$ the translation $T_a\colon x \mapsto x + a$
  is a homeomorphism of $X$ onto itself, and for each nonzero
  scalar $c$ the map $M_c\colon x \mapsto cx$ is a homeomorphism as well.
\end{theorem}
\begin{proof}
  \leanok
  $T_a$ is continuous as the sum of a constant and the identity on a
  topological additive group; its inverse is $T_{-a}$, continuous by the
  same argument. $M_c$ is continuous by \texttt{ContinuousSMul}; its
  inverse is $M_{c^{-1}}$, and
  $M_{c^{-1}}\circ M_c = \mathrm{id}$ follows from $c^{-1}(c x) = x$.
\end{proof}

\begin{theorem}[Rudin 1.22, Heine--Borel in finite dimensions]
  \label{thm:heine-borel}
  \lean{Rudin.Ch01.heine_borel}
  \leanok
  In a finite-dimensional normed space over $\mathbb{R}$ or $\mathbb{C}$,
  every closed and bounded set is compact.
\end{theorem}
\begin{proof}
  \leanok
  A finite-dimensional normed space over a locally compact field is a
  proper metric space, i.e.\ every closed ball $\overline{B}(0, R)$ is
  compact. A bounded set $s$ is contained in some such $\overline{B}(0, R)$;
  being a closed subset of a compact set, $s$ is itself compact.
\end{proof}

\begin{theorem}[corollary of 1.6, translation preserves openness]
  \label{thm:translation-preserves-open}
  \uses{thm:translation-homeo}
  \lean{Rudin.Ch01.translation_image_isOpen, Rudin.Ch01.smul_image_isOpen}
  \leanok
  Left-translation by $a$ and scaling by a nonzero $c$ send open sets
  to open sets. \emph{Proved in-project from Rudin 1.6's
  homeomorphism packaging.}
\end{theorem}
\begin{proof}
  \leanok
  Write $(a + \cdot)(s) = T_{-a}^{-1}(s)$: $y \in (a + \cdot)(s)$ iff
  $y - a \in s$ iff $T_{-a}(y) \in s$. The map $T_{-a}$ is continuous by
  Theorem~\ref{thm:translation-homeo}, so the preimage of the open $s$ is
  open. The scalar case is identical with $T_{-a}$ replaced by
  $M_{c^{-1}}$.
\end{proof}

\begin{definition}[Rudin 1.33, Minkowski functional]
  \label{def:minkowski}
  \lean{Rudin.Ch01.minkowski}
  \leanok
  For a set $s \subseteq E$ in a real topological vector space, the
  Minkowski functional of $s$ is
  $\mu_s(x) = \inf\{ t > 0 : x \in t s \}$.
\end{definition}

\begin{theorem}[Rudin 1.35, subadditivity of Minkowski functional]
  \label{thm:minkowski-subadditive}
  \uses{def:minkowski}
  \lean{Rudin.Ch01.minkowski_add_le}
  \leanok
  If $s$ is convex and absorbent, then $\mu_s(x + y) \le \mu_s(x) + \mu_s(y)$.
\end{theorem}
\begin{proof}
  \leanok
  Fix $\varepsilon > 0$. By absorbence and the definition of the gauge,
  pick $a < \mu_s(x) + \varepsilon/2$ and $b < \mu_s(y) + \varepsilon/2$
  with $u, v \in s$ and $x = a u$, $y = b v$. Convexity of $s$ gives
  $\tfrac{a}{a+b} u + \tfrac{b}{a+b} v \in s$, so
  $a u + b v = (a + b)\bigl(\tfrac{a}{a+b} u + \tfrac{b}{a+b} v\bigr) \in
  (a + b) s$. Hence $\mu_s(x + y) \le a + b < \mu_s(x) + \mu_s(y) +
  \varepsilon$. Let $\varepsilon \to 0$.
\end{proof}

\begin{theorem}[Rudin 1.18, continuity via boundedness]
  \label{thm:linear-continuity-bound}
  \lean{Rudin.Ch01.linear_continuous_of_bounded}
  \leanok
  A linear map between normed spaces is continuous iff there exists a
  constant $C$ with $\|f(x)\| \le C\|x\|$ for all $x$.
\end{theorem}
\begin{proof}
  \leanok
  Given the bound, linearity yields $\|f(x) - f(y)\| = \|f(x - y)\| \le
  C\|x - y\|$, so $f$ is Lipschitz with constant $C$, hence continuous.
  (Conversely, continuity at $0$ of a linear map forces such a bound by
  absorbing a unit ball into a chosen $0$-neighbourhood.) The Lean
  construction packages this via
  \texttt{AddMonoidHomClass.continuous\_of\_bound}.
\end{proof}

\begin{definition}[Rudin 1.27, balanced hull]
  \label{def:balanced-hull}
  \lean{Rudin.Ch01.balanced_hull}
  \leanok
  The \emph{balanced hull} of $s \subseteq E$ is the smallest balanced
  set containing $s$.
\end{definition}

\begin{theorem}[Rudin 1.28, balanced hull is balanced]
  \label{thm:balanced-hull-balanced}
  \uses{def:balanced-hull}
  \lean{Rudin.Ch01.balanced_hull_is_balanced, Rudin.Ch01.subset_balanced_hull}
  \leanok
  The balanced hull is balanced and contains the original set.
\end{theorem}
\begin{proof}
  \leanok
  Realise the hull as $\bigcup_{\|r\| \le 1} r \cdot s$. For containment
  take $r = 1$: $x = 1 \cdot x \in 1 \cdot s$ for every $x \in s$. For
  balance, take $\|a\| \le 1$ and $x \in a \cdot \bigcup_{\|r\|\le 1}
  r \cdot s$. Then $x = (a r) \cdot z$ for some $z \in s$ with
  $\|r\| \le 1$, and $\|a r\| = \|a\|\|r\| \le 1$, so $x$ lies in the
  hull.
\end{proof}

\begin{theorem}[Rudin 1.12, TVS is Hausdorff]
  \label{thm:tvs-hausdorff}
  \lean{Rudin.Ch01.t2_of_t1}
  \leanok
  Every topological vector space (in Rudin's T$_1$ sense) is Hausdorff.
\end{theorem}
\begin{proof}
  \leanok
  In any topological additive group, $T_2$ is equivalent to the
  closedness of $\{0\}$: given distinct $x \neq y$, by continuity of
  subtraction an open set separating $x - y$ from $0$ pulls back to
  disjoint open neighbourhoods of $x$ and $y$. Since $\{0\}$ is closed
  by $T_1$, the space is $T_2$.
\end{proof}

\begin{theorem}[Rudin 1.13 (d,e,f), closure properties]
  \label{thm:closure-properties}
  \lean{Rudin.Ch01.convex_closure, Rudin.Ch01.balanced_closure, Rudin.Ch01.vonNBounded_closure}
  \leanok
  In a TVS: (d) the closure of a convex set is convex; (e) the closure of
  a balanced set is balanced; (f) in a $T_1$ regular TVS, the closure of
  a bounded set is bounded.
\end{theorem}
\begin{proof}
  \leanok
  (d) For $x, y \in \overline{C}$ and $t \in [0, 1]$, the continuous map
  $(u, v) \mapsto (1 - t) u + t v$ sends $\overline{C} \times \overline{C}$
  into $\overline{\{(1 - t) u + t v : u, v \in C\}} \subseteq \overline{C}$.
  (e) The same argument with $(u) \mapsto a u$ for $\|a\| \le 1$.
  (f) Given a $0$-neighbourhood $U$, regularity yields a closed
  $0$-neighbourhood $V$ with $V \subseteq U$; boundedness of $A$ gives
  $t$ with $A \subseteq t V$; since $t V$ is closed, $\overline{A}
  \subseteq t V \subseteq t U$.
\end{proof}

\begin{theorem}[Rudin 1.14(a), balanced neighbourhood exists]
  \label{thm:balanced-nhd}
  \lean{Rudin.Ch01.nhds_zero_has_basis_balanced, Rudin.Ch01.exists_balanced_nhds_subset}
  \leanok
  The neighbourhood filter of $0$ has a basis of balanced sets;
  equivalently, every neighbourhood of $0$ contains a balanced
  neighbourhood of $0$.
\end{theorem}
\begin{proof}
  \leanok
  Given a $0$-neighbourhood $U$, define the \emph{balanced core}
  $\operatorname{bc}(U) = \bigcap_{\|a\| \ge 1} a \cdot U$. It is
  balanced by construction. Continuity of scalar multiplication at
  $(0, 0)$ yields $\delta > 0$ and a $0$-neighbourhood $W$ with $a W
  \subseteq U$ for $\|a\| \le \delta$, so $W \subseteq \bigcap_{\|a\|
  \le \delta} a^{-1} U = \operatorname{bc}(U)$ (after rescaling), making
  $\operatorname{bc}(U)$ a neighbourhood of $0$ contained in $U$.
\end{proof}

\begin{theorem}[Rudin 1.15, absorbent and bounded]
  \label{thm:absorbent-compact-bounded}
  \lean{Rudin.Ch01.absorbent_nhds_zero', Rudin.Ch01.isVonNBounded_of_isCompact}
  \leanok
  (a) Every neighbourhood of $0$ is absorbent. (b) Every compact subset
  of a TVS is bounded.
\end{theorem}
\begin{proof}
  \leanok
  (a) Fix $x \in E$ and a $0$-neighbourhood $V$. The map $a \mapsto a x$
  is continuous with $0 \cdot x = 0 \in V$, so $\{a : a x \in V\}$ is a
  neighbourhood of $0 \in \mathbb{K}$ and contains some nonzero $a$;
  then $x \in a^{-1} V$.
  (b) A compact subset of a uniform space is totally bounded. In a TVS,
  a totally bounded set is von Neumann bounded: given a $0$-neighbourhood
  $V$, pick a balanced $W$ with $W + W \subseteq V$; a finite cover of
  $K$ by translates $x_i + W$ lets one choose $t$ large enough that each
  $x_i \in t W$, hence $K \subseteq t W + W \subseteq t V$.
\end{proof}

\begin{theorem}[Rudin 1.17, continuity at zero]
  \label{thm:linear-cont-at-zero}
  \lean{Rudin.Ch01.continuous_of_linear_continuousAt_zero}
  \leanok
  A linear map between topological vector spaces that is continuous at
  $0$ is continuous everywhere.
\end{theorem}
\begin{proof}
  \leanok
  For any $x_0 \in E$, write $\Lambda = T_{\Lambda x_0} \circ \Lambda
  \circ T_{-x_0}$ (by additivity of $\Lambda$). Both translations are
  continuous (Theorem~\ref{thm:translation-homeo}), and $\Lambda$ is
  continuous at $0 = T_{-x_0}(x_0)$, so the composition is continuous at
  $x_0$.
\end{proof}

\begin{theorem}[Rudin 1.18, continuity via closed kernel]
  \label{thm:linear-cont-iff-closed-ker}
  \lean{Rudin.Ch01.linearFunctional_continuous_iff_ker_isClosed, Rudin.Ch01.linearFunctional_continuous_of_nonzero_on_open}
  \leanok
  A linear functional is continuous iff its kernel is closed. Also, a
  linear functional nonzero on some nonempty open set is automatically
  continuous.
\end{theorem}
\begin{proof}
  \leanok
  ($\Rightarrow$) $\ker \Lambda = \Lambda^{-1}\{0\}$ is the preimage of a
  closed singleton, hence closed.
  ($\Leftarrow$) Assume $\Lambda \neq 0$ and $\ker \Lambda$ closed. Pick
  $x_0$ with $\Lambda x_0 = 1$. The complement of $\ker \Lambda$ is open
  and contains $x_0$, so it contains a balanced neighbourhood $V$ of the
  form $x_0 + U$ with $U$ balanced (Theorem~\ref{thm:balanced-nhd}).
  Then $\Lambda(U)$ is a balanced subset of $\mathbb{K}$ not containing
  $-1$, which forces $|\Lambda(U)| < 1$, giving continuity at $0$;
  Theorem~\ref{thm:linear-cont-at-zero} finishes.
  For the ``nonzero on an open set'' version: if $\Lambda$ is nonzero on
  an open $s$, then $\ker \Lambda$ misses $s$, but $\ker \Lambda$ is
  either $E$ or a hyperplane; hyperplanes with nonempty open complement
  are closed, reducing to the above.
\end{proof}

\begin{theorem}[Rudin 1.20, linear maps from finite-dimensional spaces]
  \label{thm:fd-linear-continuous}
  \uses{thm:linear-cont-at-zero, thm:absorbent-compact-bounded}
  \lean{Rudin.Ch01.linear_from_finiteDimensional_continuous}
  \leanok
  Every linear map from a finite-dimensional Hausdorff topological
  vector space into a TVS is continuous.
\end{theorem}
\begin{proof}
  \leanok
  Pick a basis $e_1, \ldots, e_n$ of $E$. Writing $x = \sum_i x_i e_i$
  gives $\Lambda x = \sum_i x_i \Lambda(e_i)$, so it suffices to check
  that each coordinate projection $\pi_i\colon E \to \mathbb{K}$ is
  continuous: the target is then a finite sum of scaled continuous maps.
  On a finite-dimensional Hausdorff TVS, the $\pi_i$ are continuous
  because the space is linearly homeomorphic to $\mathbb{K}^n$ (a fact
  equivalent to the $\mathbb{K}$-module topology being uniquely
  determined).
\end{proof}

\begin{theorem}[Rudin 1.21(b), finite-dimensional subspaces are closed]
  \label{thm:fd-subspace-closed}
  \uses{thm:fd-linear-continuous}
  \lean{Rudin.Ch01.finiteDimensional_subspace_isClosed}
  \leanok
  Every finite-dimensional subspace of a Hausdorff TVS is closed.
\end{theorem}
\begin{proof}
  \leanok
  Let $F$ be a finite-dimensional subspace. The identity $F \to E$ is
  continuous, and by Theorem~\ref{thm:fd-linear-continuous} so is every
  linear bijection $\mathbb{K}^n \to F$ together with its inverse, so
  $F$ inherits a complete normed-space topology from $\mathbb{K}^n$.
  Completeness in a Hausdorff uniform space implies closedness in the
  ambient TVS.
\end{proof}

\begin{theorem}[Rudin 1.22, Riesz: locally compact TVS is finite-dim]
  \label{thm:riesz-locallycompact}
  \uses{thm:absorbent-compact-bounded, thm:balanced-nhd}
  \lean{Rudin.Ch01.finiteDimensional_of_locallyCompactSpace, Rudin.Ch01.finiteDimensional_of_totallyBounded_nhds_zero}
  \leanok
  Every locally compact topological vector space has finite dimension.
\end{theorem}
\begin{proof}
  \leanok
  Let $K$ be a compact neighbourhood of $0$; $K$ is totally bounded as
  a compact subset of a uniform space. Cover $K$ by finitely many
  translates $x_1 + \tfrac{1}{2} K, \ldots, x_m + \tfrac{1}{2} K$ and
  let $F = \operatorname{span}\{x_1, \ldots, x_m\}$. Iterating, $K
  \subseteq F + 2^{-n} K$ for every $n$; since $K$ is bounded and $F$
  is closed by Theorem~\ref{thm:fd-subspace-closed}, intersecting gives
  $K \subseteq F$. But $K$ absorbs $E$
  (Theorem~\ref{thm:absorbent-compact-bounded}(a)), so $E \subseteq F$,
  i.e.\ $\dim E \le m$.
\end{proof}

\begin{theorem}[Rudin 1.28(a), iterated translation bound]
  \label{thm:nsmul-norm}
  \lean{Rudin.Ch01.norm_nsmul_le_nat}
  \leanok
  In a seminormed additive commutative group,
  $\|n \cdot x\| \le n \cdot \|x\|$.
\end{theorem}
\begin{proof}
  \leanok
  Immediate by induction on $n$ using the triangle inequality:
  $\|(n+1) x\| = \|n x + x\| \le \|n x\| + \|x\| \le n\|x\| + \|x\| =
  (n+1)\|x\|$, with base case $\|0\| = 0$.
\end{proof}

% --- Pending: signature in Lean (Rudin/Pending.lean), proof not formalised ---

\begin{theorem}[Rudin 1.10, separation of compact and closed]
  \label{thm:tvs-separation}
  \lean{Rudin.Pending.separation_compact_closed}
  \leanok
  Let $X$ be a TVS, $K \subseteq X$ compact, $C \subseteq X$ closed,
  $K \cap C = \emptyset$. Then there exists $V \in \mathcal N(0)$ with
  $(K + V) \cap (C + V) = \emptyset$.
\end{theorem}

\begin{theorem}[Rudin 1.11, regularity of TVS]
  \label{thm:tvs-regular}
  \uses{thm:tvs-separation}
  \lean{Rudin.Pending.exists_closure_subset_of_nhds_zero}
  \leanok
  Every neighbourhood of $0$ in a TVS contains the closure of some
  neighbourhood of $0$ (the special case $K = \{0\}$ of the previous
  theorem).
\end{theorem}

\begin{theorem}[Rudin 1.13(a), closure as intersection]
  \label{thm:closure-eq-iInter-add}
  \lean{Rudin.Pending.closure_eq_iInter_add_nhds_zero}
  \leanok
  $\overline{A} = \bigcap_{V \in \mathcal N(0)} (A + V)$.
\end{theorem}

\begin{theorem}[Rudin 1.13(b), closure of sums]
  \label{thm:closure-add-subset}
  \lean{Rudin.Pending.closure_add_closure_subset_closure_add}
  \leanok
  $\overline{A} + \overline{B} \subseteq \overline{A + B}$.
\end{theorem}

\begin{definition}[Rudin 1.13(c), closure of a subspace]
  \label{def:submodule-closure}
  \lean{Rudin.Pending.submoduleClosure}
  \leanok
  The closure of a linear subspace is again a linear subspace; mathlib
  packages this as \texttt{Submodule.topologicalClosure}.
\end{definition}

\begin{theorem}[Rudin 1.14(b), convex balanced neighbourhood]
  \label{thm:convex-balanced-nhd}
  \uses{thm:balanced-nhd}
  \lean{Rudin.Pending.exists_convex_balanced_nhds_subset}
  \leanok
  In a real locally convex TVS, every convex neighbourhood of $0$
  contains a convex \emph{and} balanced neighbourhood of $0$.
\end{theorem}

\begin{theorem}[Rudin 1.15(c), bounded neighbourhood scaling base]
  \label{thm:nhds-basis-smul-bounded}
  \lean{Rudin.Pending.nhds_basis_smul_of_bounded}
  \leanok
  If $V$ is a bounded neighbourhood of $0$ and $\delta_n \to 0$
  (nonzero), then $\{\delta_n V\}$ is a countable local base at $0$.
\end{theorem}

\begin{theorem}[Rudin 1.21(a), $\mathbb K^n$-isomorphism is a homeomorphism]
  \label{thm:fd-iso-homeomorph}
  \uses{thm:fd-linear-continuous}
  \lean{Rudin.Pending.linearEquiv_finrank_isHomeomorph}
  \leanok
  Every linear isomorphism $\mathbb K^n \to Y$ onto an
  $n$-dimensional Hausdorff TVS is a homeomorphism.
\end{theorem}

\begin{theorem}[Rudin 1.23, locally bounded + Heine--Borel implies finite-dim]
  \label{thm:locallyBounded-HB-finiteDim}
  \uses{thm:closure-properties, thm:riesz-locallycompact}
  \lean{Rudin.Pending.finiteDimensional_of_locallyBounded_heineBorel}
  \leanok
  A locally bounded TVS with the Heine--Borel property is
  finite-dimensional.
\end{theorem}

\begin{theorem}[Rudin 1.24, countable base implies metrisable]
  \label{thm:countable-base-metrisable}
  \uses{thm:balanced-nhd}
  \lean{Rudin.Pending.exists_invariant_metric_of_countable_basis}
  \leanok
  A TVS with a countable local base at $0$ admits a compatible
  translation-invariant pseudo-metric whose open balls at $0$ are
  balanced (and may also be chosen convex if the space is locally
  convex).
\end{theorem}

\begin{theorem}[Rudin 1.26, dilation principle]
  \label{thm:dilation-principle}
  \lean{Rudin.Pending.dilation_principle}
  \leanok
  Suppose $(X, d_1)$ is complete, $E \subseteq X$ closed, and
  $f \colon E \to Y$ continuous with $d_2(f(x'), f(x'')) \ge d_1(x',
  x'')$ for all $x', x'' \in E$. Then $f(E)$ is closed in $Y$.
\end{theorem}

\begin{theorem}[Rudin 1.27, F-subspace is closed]
  \label{thm:fSpace-subspace-closed}
  \uses{thm:dilation-principle}
  \lean{Rudin.Pending.fSpace_subspace_isClosed}
  \leanok
  Every subspace $Y$ of a TVS $X$ that is itself an F-space (in the
  inherited topology) is closed in $X$.
\end{theorem}

\begin{theorem}[Rudin 1.28(b), rescaling a null sequence]
  \label{thm:null-rescale}
  \uses{thm:nsmul-norm}
  \lean{Rudin.Pending.exists_smul_tendsto_zero_of_tendsto_zero}
  \leanok
  If $x_n \to 0$ in a metrisable TVS, there exist positive scalars
  $\gamma_n \to \infty$ such that $\gamma_n x_n \to 0$.
\end{theorem}

\begin{theorem}[Rudin 1.30, sequential characterisation of bounded sets]
  \label{thm:bounded-iff-smul-tendsto-zero}
  \lean{Rudin.Pending.isVonNBounded_iff_smul_tendsto_zero}
  \leanok
  $E \subseteq X$ is bounded iff for every sequence $(x_n)$ in $E$ and
  every scalar sequence $\alpha_n \to 0$, one has $\alpha_n x_n \to 0$.
\end{theorem}

\begin{theorem}[Rudin 1.32, continuous iff bounded for linear maps]
  \label{thm:linear-cont-iff-bounded}
  \uses{thm:bounded-iff-smul-tendsto-zero, thm:null-rescale}
  \lean{Rudin.Pending.linear_continuous_iff_mapsBounded}
  \leanok
  For a linear map between TVSs (with the source metrisable),
  continuity is equivalent to mapping bounded sets to bounded sets.
\end{theorem}

\begin{theorem}[Rudin 1.36, Minkowski functionals are separating seminorms]
  \label{thm:minkowski-separating-seminorms}
  \uses{def:minkowski, thm:minkowski-subadditive}
  \lean{Rudin.Pending.minkowski_separating_seminorms}
  \leanok
  Given a convex balanced local base $\mathcal B$ of a real TVS, the
  family of Minkowski functionals $\{\mu_V : V \in \mathcal B\}$ is a
  separating family of continuous seminorms.
\end{theorem}

\begin{theorem}[Rudin 1.37, separating seminorms induce a topology]
  \label{thm:topology-of-seminorms}
  \uses{thm:minkowski-separating-seminorms}
  \lean{Rudin.Pending.topology_of_separating_seminorms}
  \leanok
  A separating family of seminorms $\mathcal P$ on a vector space
  induces a unique locally convex Hausdorff topology in which each
  $p \in \mathcal P$ is continuous. (Mathlib: \texttt{WithSeminorms}.)
\end{theorem}

\begin{theorem}[Rudin 1.39, normability criterion]
  \label{thm:normable-iff-convex-bounded-nhd}
  \uses{thm:convex-balanced-nhd, def:minkowski}
  \lean{Rudin.Pending.normable_iff_convex_bounded_nhds_zero}
  \leanok
  A real TVS is normable iff $0$ admits a convex bounded neighbourhood.
\end{theorem}

\begin{theorem}[Rudin 1.42, sum of closed and finite-dimensional is closed]
  \label{thm:add-closed-fd-closed}
  \uses{thm:fd-subspace-closed}
  \lean{Rudin.Pending.add_isClosed_of_finiteDimensional}
  \leanok
  If $N$ is a closed subspace and $F$ is finite-dimensional, then
  $N + F$ is closed.
\end{theorem}

\begin{theorem}[Rudin 1.44, $C(\Omega)$ is a Fréchet space]
  \label{thm:CO-frechet}
  \lean{Rudin.Ch01.continuousFunctions_isFrechetSpace}
  \leanok
  $C(\Omega, E)$, with the topology of uniform convergence on a chosen
  exhaustion of $\Omega$ by compact sets, is a Fréchet space whenever
  $\Omega$ is weakly locally compact and $\sigma$-compact and $E$ is a
  complete normed real vector space (e.g.\ $\Omega$ open in $\mathbb
  R^n$, $E = \mathbb R$ or $\mathbb C$).
\end{theorem}
\begin{proof}
  \leanok
  Each Fréchet-space ingredient is a mathlib instance on $C(\Omega,
  E)$ with the compact-open topology (= uniform convergence on
  compacta):
  topological additive group and continuous scalar multiplication
  come from the algebra/module structure on \texttt{ContinuousMap};
  local convexity is \texttt{ContinuousMap.instLocallyConvexSpace};
  countable generation of the uniformity follows from
  $\Omega$ being weakly locally compact and $\sigma$-compact together
  with the countably-generated uniformity of $E$;
  completeness is \texttt{ContinuousMap.instCompleteSpaceOfCompactly\-CoherentSpace}
  (weakly locally compact $\Rightarrow$ compactly coherent). The
  Lean proof is just a five-fold \texttt{inferInstance}.
\end{proof}

\begin{theorem}[Rudin 1.45, $H(\Omega)$ is Fréchet with Heine--Borel]
  \label{thm:H-omega-frechet-HB}
  \uses{thm:CO-frechet}
  \lean{Rudin.Pending.holomorphic_isFrechetSpace_heineBorel}
  $H(\Omega)$ (holomorphic functions on open $\Omega \subseteq
  \mathbb C$) is a Fréchet space with the Heine--Borel property under
  the compact-open topology. \emph{Pending: the Heine--Borel half is
  Montel's theorem, which is not yet in mathlib.}
\end{theorem}

\begin{theorem}[Rudin 1.45 (Fréchet part, $\Omega = \mathbb C$):
  entire functions are a closed subspace of $C(\mathbb C, \mathbb C)$]
  \label{thm:entire-isClosed}
  \uses{thm:CO-frechet}
  \lean{Rudin.Ch01.entire_isClosed, Rudin.Ch01.entireSubmodule}
  \leanok
  The set of entire functions
  $\{f \in C(\mathbb C, \mathbb C) : f \text{ is differentiable on }
  \mathbb C\}$ is a closed $\mathbb C$-submodule of
  $C(\mathbb C, \mathbb C)$. Together with
  Theorem~\ref{thm:CO-frechet} for $\Omega = \mathbb C$, this shows
  that the entire functions, with the compact-open topology, form a
  Fréchet space (closed subspace of a Fréchet space).
\end{theorem}
\begin{proof}
  \leanok
  A sequence $F_n \to F$ in $C(\mathbb C, \mathbb C)$ converges in the
  compact-open topology iff locally uniformly on $\mathbb C$
  (\texttt{ContinuousMap.tendsto\_iff\_tendsto\-LocallyUniformly};
  $\mathbb C$ is weakly locally compact). Mathlib's
  \texttt{TendstoLocallyUniformlyOn.differentiableOn} (with $U$ taken
  to be $\mathbb C$ itself) then shows that the limit of locally
  uniformly convergent differentiable functions is differentiable.
  Closure of $\{f : \text{Differentiable}\}$ under cluster points
  follows. The submodule structure (closure under sum, scalar
  multiplication, and $0$) is immediate from
  \texttt{Differentiable.add}, \texttt{Differentiable.const\_smul},
  and \texttt{differentiable\_const}.
\end{proof}

\begin{theorem}[Rudin 1.46, $C^\infty(\Omega)$ and $\mathcal D_K$ are Fréchet]
  \label{thm:Cinfty-frechet}
  \uses{thm:CO-frechet}
  \lean{Rudin.Pending.smoothFunctions_isFrechetSpace}
  $C^\infty(\Omega)$ and the test-function space $\mathcal D_K(\Omega)$
  are Fréchet spaces under the seminorms $p_n(f) = \sup_{|\alpha| \le
  n}\sup_{x \in K_n} |\partial^\alpha f(x)|$.
  \emph{Pending: completeness of $\mathcal D^n_K(E, F)$ is not yet
  available in mathlib (only the TVS / locally convex part is).}
\end{theorem}

\begin{theorem}[Rudin 1.46 (TVS + locally convex part):
  $\mathcal D^n_K(E, F)$ is a locally convex topological vector space]
  \label{thm:DKn-LCS}
  \lean{Rudin.Ch01.testFunctions_isTVS_locallyConvex}
  \leanok
  Let $E, F$ be normed real vector spaces with $F$ a normed
  $\mathbb K$-module ($\mathbb K$ a non-trivial normed field). For
  every $n \in \mathbb N \cup \{\infty\}$ and every compact
  $K \subseteq E$, the space
  $\mathcal D^{n}_{K}(E, F)$ of $n$-times continuously differentiable
  functions $E \to F$ supported in $K$ is a topological additive
  group, has continuous scalar multiplication by $\mathbb K$, and is
  locally convex over $\mathbb R$.
\end{theorem}
\begin{proof}
  \leanok
  Each conclusion is an existing mathlib instance on
  \texttt{ContDiffMapSupportedIn}:
  \texttt{isTopologicalAddGroup}, \texttt{continuousSMul}, and
  \texttt{locallyConvexSpace}. The Lean proof is a triple
  \texttt{inferInstance}.
\end{proof}

\chapter{Completeness}

\begin{theorem}[Rudin 2.2, Baire category]
  \label{thm:baire-category}
  \lean{Rudin.Ch02.baire_category}
  \leanok
  In a Baire space (in particular, any complete metric space or locally
  compact Hausdorff space), the intersection of countably many dense open
  sets is dense.
\end{theorem}
\begin{proof}
  \leanok
\end{proof}

\begin{theorem}[Rudin 2.5, Banach--Steinhaus]
  \label{thm:banach-steinhaus}
  \uses{thm:baire-category}
  \lean{Rudin.Ch02.banach_steinhaus}
  \leanok
  Let $E$ be a Banach space, $F$ a normed space, and $(g_i)_{i\in I}$
  a family of continuous linear maps $E \to F$. If the family is
  pointwise bounded, then it is uniformly bounded in operator norm.
\end{theorem}
\begin{proof}
  \leanok
\end{proof}

\begin{theorem}[Rudin 2.6, pointwise convergent $\Rightarrow$ uniformly bounded]
  \label{thm:banach-steinhaus-pointwise}
  \uses{thm:banach-steinhaus}
  \lean{Rudin.Ch02.banach_steinhaus_of_pointwise_tendsto}
  \leanok
  If a sequence $(g_n)$ of continuous linear maps from a Banach space
  converges pointwise, then the operator norms $\|g_n\|$ are uniformly
  bounded. \emph{Proved in-project: convergence $\Rightarrow$ Cauchy
  $\Rightarrow$ bounded range, then apply Banach--Steinhaus.}
\end{theorem}
\begin{proof}
  \leanok
\end{proof}

\begin{theorem}[Rudin 2.11, Open mapping]
  \label{thm:open-mapping}
  \uses{thm:baire-category}
  \lean{Rudin.Ch02.open_mapping}
  \leanok
  A surjective continuous linear map between Banach spaces is open.
\end{theorem}
\begin{proof}
  \leanok
\end{proof}

\begin{theorem}[Rudin 2.15, Closed graph]
  \label{thm:closed-graph}
  \uses{thm:open-mapping}
  \lean{Rudin.Ch02.closed_graph}
  \leanok
  A linear map between Banach spaces whose graph is closed is continuous.
\end{theorem}
\begin{proof}
  \leanok
\end{proof}

\begin{theorem}[Rudin, Riesz's lemma]
  \label{thm:riesz-lemma}
  \lean{Rudin.Ch02.riesz_lemma}
  \leanok
  For any closed proper subspace $F$ of a normed space $E$ and any
  $r < 1$, there exists $x_0 \notin F$ with $r\,\|x_0\| \le \|x_0 - y\|$
  for all $y \in F$.
\end{theorem}
\begin{proof}
  \leanok
\end{proof}

\begin{theorem}[Rudin 2.12, Bounded inverse]
  \label{thm:bounded-inverse}
  \uses{thm:open-mapping}
  \lean{Rudin.Ch02.bounded_inverse}
  \leanok
  A bijective continuous linear map between Banach spaces has a
  continuous inverse (so it is a topological isomorphism).
\end{theorem}
\begin{proof}
  \leanok
\end{proof}

\begin{definition}[meagre / first-category sets]
  \label{def:meagre}
  \lean{Rudin.Ch02.isMeagre}
  \leanok
  A set is \emph{meagre} (a.k.a.\ of the \emph{first category}) if it
  is a countable union of nowhere-dense sets.
\end{definition}

\begin{theorem}[meagre complement is dense in Baire spaces]
  \label{thm:meagre-compl-dense}
  \uses{def:meagre, thm:baire-category}
  \lean{Rudin.Ch02.meagre_compl_dense}
  \leanok
  In a Baire space, every meagre set has dense complement.
\end{theorem}
\begin{proof}
  \leanok
\end{proof}

% --- Ch02 Pending ---

\begin{theorem}[Rudin 2.4, equicontinuous + bounded $\Rightarrow$ uniform bound]
  \label{thm:equicontinuous-uniform-bound}
  \lean{Rudin.Pending.equicontinuous_uniformly_bounded}
  \leanok
  An equicontinuous family of linear maps takes every bounded subset
  of the source into a single bounded subset of the target.
\end{theorem}

\begin{theorem}[Rudin 2.7, Cauchy on second-category $\Rightarrow$ all]
  \label{thm:cauchy-second-category}
  \uses{thm:banach-steinhaus}
  \lean{Rudin.Pending.cauchy_set_eq_univ_of_second_category}
  \leanok
  If $(\Lambda_n)$ is a sequence of continuous linear maps and the set
  of $x$ with $(\Lambda_n x)$ Cauchy is of second category, then every
  $\Lambda_n x$ sequence is Cauchy.
\end{theorem}

\begin{theorem}[Rudin 2.9, compact-convex Banach--Steinhaus]
  \label{thm:banach-steinhaus-compact-convex}
  \uses{thm:baire-category}
  \lean{Rudin.Pending.banach_steinhaus_compact_convex}
  \leanok
  If a family of continuous linear maps is pointwise bounded on a
  compact convex set $K$, it is uniformly bounded on $K$.
\end{theorem}

\begin{theorem}[Rudin 2.12(c), equivalent norms via open mapping]
  \label{thm:open-mapping-norm-equiv}
  \uses{thm:open-mapping}
  \lean{Rudin.Pending.norm_equiv_of_continuousLinearEquiv}
  \leanok
  A continuous linear bijection $A : X \to Y$ between Banach spaces
  satisfies $a\|x\| \le \|A x\| \le b\|x\|$ for some $a, b > 0$.
\end{theorem}

\begin{theorem}[Rudin 2.12(d), comparable F-space topologies coincide]
  \label{thm:fSpace-topologies-eq}
  \uses{thm:open-mapping}
  \lean{Rudin.Pending.fSpace_topologies_eq_of_le}
  \leanok
  Two comparable F-space topologies on the same vector space are
  equal.
\end{theorem}

\begin{theorem}[Rudin 2.14, graph closed if continuous (Hausdorff target)]
  \label{thm:graph-closed-of-continuous}
  \lean{Rudin.Pending.graph_isClosed_of_continuous_of_t2}
  \leanok
  If $Y$ is Hausdorff and $f : X \to Y$ is continuous, the graph of
  $f$ is closed in $X \times Y$.
\end{theorem}

\begin{theorem}[Rudin 2.17, separately continuous bilinear is jointly continuous]
  \label{thm:bilinear-cont-of-separately}
  \uses{thm:banach-steinhaus}
  \lean{Rudin.Pending.bilinear_continuous_of_separately}
  \leanok
  A separately continuous bilinear map $B : X \times Y \to Z$ with
  $X$ an F-space and $Y$ metrisable is jointly continuous.
\end{theorem}

\chapter{Convexity}

\begin{theorem}[Rudin 3.2, Hahn--Banach, analytic]
  \label{thm:hahn-banach-analytic}
  \lean{Rudin.Ch03.hahn_banach_extension}
  \leanok
  Let $E$ be a real vector space, $N\colon E \to \mathbb{R}$ a sublinear
  functional, and $f$ a linear functional defined on a subspace of $E$
  with $f \le N$ on its domain. Then $f$ extends to a linear functional
  $g\colon E \to \mathbb{R}$ with $g \le N$ everywhere.
\end{theorem}
\begin{proof}
  \leanok
\end{proof}

\begin{theorem}[Rudin 3.4, Geometric Hahn--Banach]
  \label{thm:hahn-banach-separation}
  \uses{thm:hahn-banach-analytic}
  \lean{Rudin.Ch03.hahn_banach_separation}
  \leanok
  In a locally convex Hausdorff topological vector space, a compact
  convex set and a disjoint closed convex set can be strictly separated
  by a continuous linear functional.
\end{theorem}
\begin{proof}
  \leanok
\end{proof}

\begin{theorem}[Rudin 3.5, closed convex vs. point]
  \label{thm:hahn-banach-closed-point}
  \uses{thm:hahn-banach-analytic}
  \lean{Rudin.Ch03.hahn_banach_closed_point}
  \leanok
  A closed convex set $s$ and a point $x \notin s$ in a locally convex
  Hausdorff TVS can be strictly separated by a continuous linear
  functional.
\end{theorem}
\begin{proof}
  \leanok
\end{proof}

\begin{theorem}[Rudin 3.5, two distinct points]
  \label{thm:hahn-banach-point-point}
  \uses{thm:hahn-banach-closed-point}
  \lean{Rudin.Ch03.hahn_banach_point_point}
  \leanok
  In a locally convex $T_1$ TVS, any two distinct points are strictly
  separated by a continuous linear functional.
\end{theorem}
\begin{proof}
  \leanok
\end{proof}

\begin{theorem}[Rudin 3.22, Krein--Milman lemma]
  \label{thm:krein-milman-lemma}
  \uses{thm:hahn-banach-point-point}
  \lean{Rudin.Ch03.krein_milman_lemma}
  \leanok
  Every nonempty compact set in a locally convex Hausdorff TVS has an
  extreme point.
\end{theorem}
\begin{proof}
  \leanok
\end{proof}

\begin{theorem}[Rudin 3.23, Krein--Milman]
  \label{thm:krein-milman}
  \uses{thm:krein-milman-lemma, thm:hahn-banach-separation}
  \lean{Rudin.Ch03.krein_milman}
  \leanok
  In a locally convex Hausdorff topological vector space, every compact
  convex set is the closed convex hull of its extreme points.
\end{theorem}
\begin{proof}
  \leanok
\end{proof}

% --- Ch03 Pending ---

\begin{theorem}[Rudin 3.3, complex Hahn--Banach]
  \label{thm:hahn-banach-complex}
  \uses{thm:hahn-banach-analytic}
  \lean{Rudin.Pending.hahn_banach_complex}
  \leanok
  A complex-linear functional on a subspace, dominated by a seminorm,
  extends to the whole space with the same dominance.
\end{theorem}

\begin{theorem}[Rudin 3.6, continuous extension Hahn--Banach]
  \label{thm:hahn-banach-cont-extension}
  \uses{thm:hahn-banach-complex}
  \lean{Rudin.Pending.hahn_banach_continuous_extension}
  \leanok
  A continuous linear functional on a subspace of a locally convex
  TVS extends to a continuous linear functional on the whole space.
\end{theorem}

\begin{theorem}[Rudin 3.7, open convex / point separation]
  \label{thm:hahn-banach-open-point}
  \uses{thm:hahn-banach-cont-extension}
  \lean{Rudin.Pending.hahn_banach_open_point}
  \leanok
  A nonempty open convex set and a point not in it can be separated
  by a continuous linear functional.
\end{theorem}

\begin{theorem}[Rudin 3.10, separating dual induces locally convex topology]
  \label{thm:topology-separating-dual}
  \lean{Rudin.Pending.topology_of_separating_dual}
  \leanok
  A separating vector space $X'$ of linear functionals on $X$
  induces a locally convex topology on $X$ whose dual is exactly
  $X'$.
\end{theorem}

\begin{theorem}[Rudin 3.12, weak closure equals original closure for convex sets]
  \label{thm:weak-closure-convex}
  \uses{thm:hahn-banach-separation}
  \lean{Rudin.Pending.weak_closure_eq_closure_of_convex}
  \leanok
  In a locally convex space, the weak closure of a convex set equals
  its original closure.
\end{theorem}

\begin{theorem}[Rudin 3.13, Mazur's lemma]
  \label{thm:mazur}
  \uses{thm:weak-closure-convex}
  \lean{Rudin.Pending.mazur_convex_combinations}
  \leanok
  In a metrisable locally convex space, weak limits of sequences are
  strong limits of convex combinations.
\end{theorem}

\begin{theorem}[Rudin 3.16, weak-* compact metrisable in separable]
  \label{thm:weakStar-metrisable}
  \lean{Rudin.Pending.weakStar_compact_metrisable_of_separable}
  \leanok
  In a separable TVS, weak-* compact subsets of the dual are
  metrisable in the weak-* topology.
\end{theorem}

\begin{theorem}[Rudin 3.17, sequential Banach--Alaoglu]
  \label{thm:banach-alaoglu-sequential}
  \uses{thm:banach-alaoglu, thm:weakStar-metrisable}
  \lean{Rudin.Pending.banach_alaoglu_sequential}
  \leanok
  In a separable TVS, the polar of any neighbourhood of $0$ is
  sequentially compact in weak-*.
\end{theorem}

\begin{theorem}[Rudin 3.18, weakly bounded $=$ bounded]
  \label{thm:weakly-bounded-iff-bounded}
  \uses{thm:banach-steinhaus}
  \lean{Rudin.Pending.weakly_bounded_iff_bounded}
  \leanok
  In a locally convex space, weakly bounded sets are originally
  bounded (and vice versa).
\end{theorem}

\begin{theorem}[Rudin 3.20, bipolar theorem]
  \label{thm:bipolar}
  \uses{thm:weak-closure-convex}
  \lean{Rudin.Pending.bipolar_eq_weak_closure}
  \leanok
  For a balanced convex $E \subseteq X$, the bipolar $E^{\circ\circ}$
  is the weak closure of $E$.
\end{theorem}

\begin{theorem}[Rudin 3.21, convex closed iff bipolar self]
  \label{thm:convex-closed-iff-bipolar}
  \uses{thm:bipolar}
  \lean{Rudin.Pending.convex_closed_iff_bipolar}
  \leanok
  A convex set is closed iff it equals its bipolar.
\end{theorem}

\begin{theorem}[Rudin 3.25, Milman]
  \label{thm:milman}
  \uses{thm:krein-milman}
  \lean{Rudin.Pending.milman}
  \leanok
  If $K$ is compact and the closed convex hull of $K$ is also compact
  in a locally convex space, then the extreme points of the closed
  convex hull lie in $K$.
\end{theorem}

\begin{theorem}[Rudin 3.27, vector-valued integral]
  \label{thm:vector-integral}
  \lean{Rudin.Pending.vector_integral_exists}
  \leanok
  For $f$ continuous from a compact Hausdorff $(Q, \mu)$ to a Fréchet
  space, there exists $y$ with $\Lambda y = \int (\Lambda \circ f)
  d\mu$ for every continuous linear $\Lambda$.
\end{theorem}

\begin{theorem}[Rudin 3.29, weak holomorphy implies strong]
  \label{thm:weak-holomorphic-strong}
  \lean{Rudin.Pending.weak_holomorphic_iff_strong}
  \leanok
  Weakly holomorphic maps into a Fréchet space are strongly
  holomorphic.
\end{theorem}

\chapter{Duality in Banach Spaces}

\begin{theorem}[Rudin 4.3, Banach--Alaoglu]
  \label{thm:banach-alaoglu}
  \lean{Rudin.Ch04.banach_alaoglu, Rudin.Ch04.banach_alaoglu_polar}
  \leanok
  The closed unit ball of the continuous dual of a normed space is
  compact in the weak-$*$ topology; more generally, the polar of any
  neighborhood of the origin is weak-$*$ compact.
\end{theorem}
\begin{proof}
  \leanok
\end{proof}

\begin{definition}[Rudin 4.4, canonical embedding into the bidual]
  \label{def:double-dual-embedding}
  \lean{Rudin.Ch04.doubleDualEmbedding}
  \leanok
  The canonical map $E \to E^{**}$, $x \mapsto (\varphi \mapsto \varphi(x))$,
  is an isometric linear embedding. $E$ is called \emph{reflexive} if
  this map is surjective.
\end{definition}

\begin{definition}[polar set]
  \label{def:polar}
  \lean{Rudin.Ch04.polar}
  \leanok
  For $s \subseteq E$, the \emph{polar} is
  $\{\varphi \in E^{*} : |\varphi(x)| \le 1 \text{ for all } x \in s\}$.
\end{definition}

\begin{theorem}[polar is closed]
  \label{thm:polar-closed}
  \uses{def:polar}
  \lean{Rudin.Ch04.polar_isClosed}
  \leanok
  The polar of any set is closed in the strong dual.
\end{theorem}
\begin{proof}
  \leanok
\end{proof}

% --- Ch04 Pending ---

\begin{theorem}[Rudin 4.1, dual norm makes $X^*$ Banach]
  \label{thm:dualNorm-banach}
  \lean{Rudin.Pending.dualNorm_isBanach}
  \leanok
  $\|\Lambda\| = \sup_{\|x\| \le 1} |\Lambda x|$ makes the dual of a
  normed space into a Banach space.
\end{theorem}

\begin{theorem}[Rudin 4.5, canonical embedding is isometric]
  \label{thm:canonical-embedding-isometry}
  \uses{def:double-dual-embedding}
  \lean{Rudin.Pending.canonical_embedding_isometry}
  \leanok
  The canonical embedding $X \to X^{**}$ is an isometry.
\end{theorem}

\begin{theorem}[Rudin 4.6, annihilator double]
  \label{thm:annihilator-double}
  \uses{thm:hahn-banach-cont-extension}
  \lean{Rudin.Pending.annihilator_double}
  \leanok
  $M^\perp$ is a closed subspace of $X^*$, and
  $M^{\perp\perp} \cap X = \overline{M}$ (under the canonical
  embedding).
\end{theorem}

\begin{theorem}[Rudin 4.7, annihilator codimension duality]
  \label{thm:annihilator-codim}
  \uses{thm:annihilator-double}
  \lean{Rudin.Pending.annihilator_codim}
  \leanok
  $(M^\perp)^\perp = \overline{M}$, and dimensions/codimensions of
  subspaces and their annihilators correspond appropriately.
\end{theorem}

\begin{theorem}[Rudin 4.9, duals of subspaces and quotients]
  \label{thm:dual-subspace-quotient}
  \uses{thm:hahn-banach-cont-extension}
  \lean{Rudin.Pending.dual_of_subspace_quotient}
  \leanok
  For a closed subspace $M \subseteq X$,
  $M^* \cong X^*/M^\perp$ and $(X/M)^* \cong M^\perp$.
\end{theorem}

\begin{theorem}[Rudin 4.10, norm-preserving Hahn--Banach]
  \label{thm:hahn-banach-norm-preserving}
  \uses{thm:hahn-banach-cont-extension}
  \lean{Rudin.Pending.hahn_banach_norm_preserving}
  \leanok
  Every continuous linear functional on a closed subspace extends to
  the whole space with the same norm.
\end{theorem}

\begin{theorem}[Rudin 4.12, adjoint continuity characterisation]
  \label{thm:adjoint-cont-iff}
  \lean{Rudin.Pending.adjoint_continuous_iff}
  \leanok
  $T$ is continuous iff its adjoint $T^*$ is weak-* continuous.
\end{theorem}

\begin{theorem}[Rudin 4.13, $T(U) \subseteq V$ iff $\|T\| \le 1$]
  \label{thm:operatorNorm-le-one}
  \lean{Rudin.Pending.operatorNorm_le_one_iff}
  \leanok
  An operator maps the open unit ball into the open unit ball iff its
  operator norm is $\le 1$.
\end{theorem}

\begin{theorem}[Rudin 4.14, openness via adjoint]
  \label{thm:open-via-adjoint}
  \uses{thm:open-mapping}
  \lean{Rudin.Pending.open_iff_adjoint_bounded_below}
  \leanok
  $T : X \to Y$ between Banach spaces is open iff $T^*$ has closed
  range and is bounded below.
\end{theorem}

\begin{theorem}[Rudin 4.18, compact operators ideal]
  \label{thm:compactOps-ideal}
  \lean{Rudin.Pending.compactOps_isClosedTwoSidedIdeal}
  \leanok
  Compact operators form a closed two-sided ideal in $\mathcal B(X)$.
\end{theorem}

\begin{theorem}[Rudin 4.19, adjoint of compact is compact]
  \label{thm:adjoint-compact}
  \lean{Rudin.Pending.adjoint_compact_of_compact}
  \leanok
  The adjoint of a compact operator is compact.
\end{theorem}

\begin{theorem}[Rudin 4.21, functional vanishing on subspace]
  \label{thm:functional-vanishing}
  \uses{thm:hahn-banach-cont-extension}
  \lean{Rudin.Pending.exists_functional_eq_one_off_subspace}
  \leanok
  For $M$ closed and $x \notin M$, there is a continuous linear
  functional vanishing on $M$ with $\Lambda x = 1$.
\end{theorem}

\begin{theorem}[Rudin 4.24, spectrum of compact operator]
  \label{thm:spectrum-compactOp}
  \uses{thm:compactOps-ideal, thm:riesz-lemma}
  \lean{Rudin.Pending.spectrum_compactOperator}
  \leanok
  Spectrum of a compact operator on a Banach space is countable, with
  $0$ the only possible limit point; nonzero spectrum consists of
  eigenvalues with finite-dimensional eigenspaces.
\end{theorem}

\begin{theorem}[Rudin 4.25, Fredholm alternative]
  \label{thm:fredholm}
  \uses{thm:spectrum-compactOp}
  \lean{Rudin.Pending.fredholm_alternative}
  \leanok
  Fredholm alternative for compact perturbations of the identity.
\end{theorem}

\chapter{Some Applications}

\begin{theorem}[Rudin 5.7, Stone--Weierstrass (real)]
  \label{thm:stone-weierstrass}
  \lean{Rudin.Ch05.stone_weierstrass}
  \leanok
  If $X$ is a compact Hausdorff space and $A$ is a real subalgebra of
  $C(X,\mathbb{R})$ that separates points, then $A$ is dense in
  $C(X,\mathbb{R})$.
\end{theorem}
\begin{proof}
  \leanok
\end{proof}

\begin{theorem}[Rudin 5.8, Stone--Weierstrass ($\mathbb{R},\mathbb{C}$)]
  \label{thm:stone-weierstrass-star}
  \lean{Rudin.Ch05.stone_weierstrass_star}
  \leanok
  For $\mathbb{K} \in \{\mathbb{R}, \mathbb{C}\}$ and $X$ compact, a
  self-adjoint star-subalgebra of $C(X, \mathbb{K})$ that separates
  points is dense.
\end{theorem}
\begin{proof}
  \leanok
\end{proof}

% --- Ch05 Pending ---

\begin{theorem}[Rudin 5.1, continuous convex bounded above on bounded sets]
  \label{thm:continuousConvex-boundedAbove}
  \lean{Rudin.Pending.continuousConvex_boundedAbove_on_bounded}
  \leanok
  A continuous convex function on a locally convex space is bounded
  above on bounded sets.
\end{theorem}

\begin{theorem}[Rudin 5.5, continuous partition of unity]
  \label{thm:partition-of-unity}
  \lean{Rudin.Pending.partition_of_unity_continuous}
  \leanok
  Existence of continuous partitions of unity subordinate to a given
  open cover (technical lemma used throughout Ch.~5).
\end{theorem}

\begin{theorem}[Rudin 5.10, Runge approximation]
  \label{thm:runge}
  \lean{Rudin.Pending.runge_approximation}
  \leanok
  Every holomorphic function on an open subset of $\mathbb C$ can be
  uniformly approximated on compacta by rational functions whose
  poles lie outside the domain.
\end{theorem}

\begin{theorem}[Rudin 5.18, Müntz--Szász]
  \label{thm:muntz-szasz}
  \lean{Rudin.Pending.muntz_szasz}
  \leanok
  The span of $\{x^{n_k}\}$ is dense in $C[0,1]$ iff
  $\sum 1/n_k = \infty$.
\end{theorem}

\begin{theorem}[Rudin 5.21, partial continuous extension]
  \label{thm:partial-cont-extension}
  \lean{Rudin.Pending.partial_continuous_extension}
  \leanok
  Hahn--Banach-style continuous extension theorem for partial
  functions.
\end{theorem}

\begin{theorem}[Rudin 5.23, Markov--Kakutani fixed point]
  \label{thm:markov-kakutani}
  \uses{thm:krein-milman}
  \lean{Rudin.Pending.markov_kakutani}
  \leanok
  A commuting family of continuous affine self-maps of a compact
  convex set has a common fixed point.
\end{theorem}

\begin{theorem}[Rudin 5.27, Kakutani--Yosida ergodic]
  \label{thm:kakutani-yosida}
  \lean{Rudin.Pending.kakutani_yosida_ergodic}
  \leanok
  Kakutani--Yosida ergodic theorem on Banach spaces.
\end{theorem}

\chapter{Test Functions and Distributions}

\begin{definition}[Schwartz space]
  \label{def:schwartz}
  \lean{Rudin.Ch06.Schwartz}
  \leanok
  The space $\mathcal{S}(E,F)$ of rapidly decreasing smooth functions
  from a finite-dimensional real vector space $E$ to a normed space
  $F$, with the Fréchet topology generated by the seminorms
  $\|x\|^k \|f^{(n)}(x)\|$.
\end{definition}

\begin{theorem}[Rudin 6.4, differentiation is continuous on test functions]
  \label{thm:fderiv-schwartz}
  \uses{def:schwartz}
  \lean{Rudin.Ch06.fderivSchwartz}
  \leanok
  The Fréchet derivative defines a continuous linear map
  $\mathcal{S}(E,F) \to \mathcal{S}(E, E\!\to_L\! F)$.
\end{theorem}
\begin{proof}
  \leanok
\end{proof}

% --- Ch06 Pending ---

\begin{theorem}[Rudin 6.3, $\mathcal D(\Omega)$ inductive limit]
  \label{thm:testFunctions-inductiveLimit}
  \uses{def:schwartz}
  \lean{Rudin.Pending.testFunctions_inductiveLimit}
  \leanok
  $\mathcal D(\Omega) = \bigcup_K \mathcal D_K(\Omega)$ carries the
  strict inductive limit topology of the $\mathcal D_K$.
\end{theorem}

\begin{theorem}[Rudin 6.5, distribution iff continuous on each $\mathcal D_K$]
  \label{thm:distribution-iff-DK}
  \uses{thm:testFunctions-inductiveLimit}
  \lean{Rudin.Pending.distribution_continuous_iff_on_DK}
  \leanok
  A linear functional on $\mathcal D(\Omega)$ is a distribution iff
  its restriction to each $\mathcal D_K$ is continuous.
\end{theorem}

\begin{theorem}[Rudin 6.6, sequential characterisation of distributions]
  \label{thm:distribution-sequential}
  \uses{thm:distribution-iff-DK}
  \lean{Rudin.Pending.distribution_iff_sequential}
  \leanok
  A linear $\Lambda$ on $\mathcal D(\Omega)$ is a distribution iff
  $\Lambda \varphi_n \to 0$ whenever $\varphi_n \to 0$ in some
  $\mathcal D_K$.
\end{theorem}

\begin{theorem}[Rudin 6.8, seminorm bound characterisation]
  \label{thm:distribution-seminorm}
  \uses{thm:distribution-iff-DK}
  \lean{Rudin.Pending.distribution_iff_seminorm_bound}
  \leanok
  $\Lambda$ is a distribution iff for every compact $K$ there are
  $c, N$ with
  $|\Lambda \varphi| \le c \sum_{|\alpha| \le N} \|\partial^\alpha
  \varphi\|_\infty$.
\end{theorem}

\begin{theorem}[Rudin 6.13, distribution derivatives]
  \label{thm:distribution-derivative}
  \lean{Rudin.Pending.distribution_derivative_exists}
  \leanok
  Every distribution has distribution derivatives of all orders,
  defined by $(\partial^\alpha T)(\varphi) = (-1)^{|\alpha|}
  T(\partial^\alpha \varphi)$.
\end{theorem}

\begin{theorem}[Rudin 6.17, $\mathcal D'(\Omega)$ is sequentially complete]
  \label{thm:distributions-seqComplete}
  \lean{Rudin.Pending.distributions_sequentiallyComplete}
  \leanok
  Pointwise limits of distributions are distributions.
\end{theorem}

\begin{theorem}[Rudin 6.20, localisation principle]
  \label{thm:distribution-local}
  \lean{Rudin.Pending.distribution_local_principle}
  \leanok
  A distribution that vanishes on a neighbourhood of every point of
  an open set vanishes on that open set.
\end{theorem}

\begin{theorem}[Rudin 6.24, point-support distributions]
  \label{thm:distribution-pointSupport}
  \lean{Rudin.Pending.distribution_pointSupport_eq_diracDerivatives}
  \leanok
  A distribution with point support is a finite linear combination
  of derivatives of the Dirac delta.
\end{theorem}

\begin{theorem}[Rudin 6.25, compactly supported distributions $= \mathcal E'$]
  \label{thm:E-prime}
  \lean{Rudin.Pending.compactlySupported_distributions_eq_dual_smooth}
  \leanok
  $\mathcal E'(\Omega) = (\mathcal E(\Omega))^*$ as the space of
  compactly supported distributions.
\end{theorem}

\begin{theorem}[Rudin 6.30, convolution of distribution and test function]
  \label{thm:convolution-T-phi}
  \lean{Rudin.Pending.convolution_distribution_test}
  \leanok
  $T \star \varphi$ is smooth, and derivatives commute with the
  convolution.
\end{theorem}

\begin{theorem}[Rudin 6.32, smooth approximation of distributions]
  \label{thm:distribution-smooth-approx}
  \uses{thm:convolution-T-phi}
  \lean{Rudin.Pending.distribution_approx_smooth}
  \leanok
  Every distribution is a limit of smooth functions in $\mathcal D'$.
\end{theorem}

\begin{theorem}[Rudin 6.36, convolution of distributions]
  \label{thm:convolution-distributions}
  \uses{thm:E-prime}
  \lean{Rudin.Pending.convolution_distributions}
  \leanok
  $T \star S$ is well-defined when at least one of $T, S$ has compact
  support.
\end{theorem}

\chapter{Fourier Transforms}

\begin{theorem}[Rudin 7.4, Fourier transform on Schwartz space]
  \label{thm:fourier-schwartz}
  \uses{def:schwartz}
  \lean{Rudin.Ch07.fourierSchwartz}
  \leanok
  The Fourier transform is a continuous linear self-map of
  $\mathcal{S}(V, E)$.
\end{theorem}
\begin{proof}
  \leanok
\end{proof}

\begin{theorem}[Rudin 7.9, Plancherel]
  \label{thm:plancherel}
  \uses{thm:fourier-schwartz}
  \lean{Rudin.Ch07.fourierL2, Rudin.Ch07.plancherel}
  \leanok
  The Fourier transform extends to a unitary automorphism of $L^2(E;F)$;
  in particular, $\|\widehat{f}\|_{L^2} = \|f\|_{L^2}$.
\end{theorem}
\begin{proof}
  \leanok
\end{proof}

\begin{theorem}[Rudin 7.7, Fourier inversion on Schwartz space]
  \label{thm:fourier-inversion}
  \uses{thm:fourier-schwartz}
  \lean{Rudin.Ch07.fourier_inversion}
  \leanok
  For $f \in \mathcal{S}(V, E)$, $\mathcal{F}^{-1} \mathcal{F} f = f$.
\end{theorem}
\begin{proof}
  \leanok
\end{proof}

% --- Ch07 Pending ---

\begin{theorem}[Rudin 7.1, Fourier transform on $L^1$]
  \label{thm:fourierTransform-L1}
  \lean{Rudin.Pending.fourierTransform_L1}
  \leanok
  $\widehat{f}(\xi) = \int f(x) e^{-2\pi i x \cdot \xi} dx$ defines a
  bounded operator $L^1 \to L^\infty$.
\end{theorem}

\begin{theorem}[Rudin 7.2, Riemann--Lebesgue]
  \label{thm:riemann-lebesgue}
  \uses{thm:fourierTransform-L1}
  \lean{Rudin.Pending.riemann_lebesgue}
  \leanok
  $\widehat{f} \in C_0(\mathbb R^n)$ for $f \in L^1$.
\end{theorem}

\begin{theorem}[Rudin 7.6, Fourier exchanges convolution and multiplication]
  \label{thm:fourier-convolution}
  \uses{thm:fourierTransform-L1}
  \lean{Rudin.Pending.fourier_convolution}
  \leanok
  $\widehat{f \star g} = \widehat{f} \cdot \widehat{g}$.
\end{theorem}

\begin{theorem}[Rudin 7.13--7.15, Paley--Wiener for distributions]
  \label{thm:paley-wiener}
  \uses{thm:E-prime}
  \lean{Rudin.Pending.paley_wiener}
  \leanok
  A tempered distribution $T$ has compact support iff $\widehat{T}$
  extends to an entire function of exponential type.
\end{theorem}

\begin{theorem}[Rudin 7.19, Schwartz invariance]
  \label{thm:fourier-schwartz-invariant}
  \uses{thm:fourier-schwartz}
  \lean{Rudin.Pending.fourier_schwartz_invariant}
  \leanok
  Schwartz class is invariant under Fourier transform; the inverse
  is $\mathcal F^{-1} f(x) = \mathcal F f(-x)$.
\end{theorem}

\begin{theorem}[Rudin 7.23, Bochner]
  \label{thm:bochner}
  \lean{Rudin.Pending.bochner}
  \leanok
  A continuous $\varphi$ on $\mathbb R^n$ is positive definite iff
  it is the Fourier transform of a nonnegative finite Borel measure.
\end{theorem}

\chapter{Applications to Differential Equations}

\begin{theorem}[Fourier transforms differentiation into multiplication]
  \label{thm:fourier-derivative}
  \uses{thm:fourier-schwartz}
  \lean{Rudin.Ch08.fourier_fderiv_eq}
  \leanok
  For $f \in \mathcal{S}(V,E)$,
  $\widehat{\partial f}(\xi) = 2\pi i\,\xi \cdot \widehat{f}(\xi)$,
  reducing constant-coefficient linear PDEs in $\mathcal{S}$ to
  algebraic (multiplication) problems.
\end{theorem}
\begin{proof}
  \leanok
\end{proof}

% --- Ch08 Pending ---

\begin{theorem}[Rudin 8.3, fundamental solution lemma]
  \label{thm:fundamental-solution}
  \lean{Rudin.Pending.fundamental_solution_exists}
  \leanok
  For polynomials $P$ of degree $N$, there is a fundamental solution
  $E$ for $P(D)$ (a distribution with $P(D) E = \delta$).
\end{theorem}

\begin{theorem}[Rudin 8.4, Malgrange--Ehrenpreis]
  \label{thm:malgrange-ehrenpreis}
  \uses{thm:fundamental-solution}
  \lean{Rudin.Pending.malgrange_ehrenpreis}
  \leanok
  Every nonzero constant-coefficient linear PDE has a fundamental
  solution.
\end{theorem}

\begin{theorem}[Rudin 8.5--8.6, elliptic regularity]
  \label{thm:elliptic-regularity}
  \uses{thm:malgrange-ehrenpreis}
  \lean{Rudin.Pending.elliptic_regularity}
  \leanok
  If $P(D) u = f$ with $f$ smooth and $P$ elliptic, then $u$ is
  smooth.
\end{theorem}

\begin{theorem}[Rudin 8.9, local solvability]
  \label{thm:local-solvability}
  \uses{thm:malgrange-ehrenpreis}
  \lean{Rudin.Pending.local_solvability}
  \leanok
  Every constant-coefficient linear PDE is locally solvable.
\end{theorem}

\begin{theorem}[Rudin 8.12, hypoelliptic operators]
  \label{thm:hypoelliptic}
  \lean{Rudin.Pending.hypoelliptic_iff}
  \leanok
  $P(D)$ is hypoelliptic iff its fundamental solution is smooth away
  from the origin.
\end{theorem}

\begin{theorem}[Rudin 8.14, quasi-elliptic implies hypoelliptic]
  \label{thm:quasiElliptic-hypoelliptic}
  \uses{thm:hypoelliptic}
  \lean{Rudin.Pending.quasiElliptic_hypoelliptic}
  \leanok
  Quasi-elliptic operators are hypoelliptic.
\end{theorem}

\chapter{Tauberian Theory}

\begin{theorem}[Abel's limit theorem]
  \label{thm:abel-limit}
  \lean{Rudin.Ch09.abel_limit}
  \leanok
  If $\sum a_n$ converges to $l$, then $\sum a_n x^n \to l$ as
  $x \to 1^-$. (The elementary abelian partner to Wiener's Tauberian
  theorem, which is not yet available in mathlib.)
\end{theorem}
\begin{proof}
  \leanok
\end{proof}

% --- Ch09 Pending ---

\begin{theorem}[Rudin 9.3, Wiener's lemma]
  \label{thm:wiener-lemma}
  \uses{thm:fourier-convolution}
  \lean{Rudin.Pending.wiener_lemma}
  \leanok
  If $\varphi \in L^1(\mathbb R^n)$ has nowhere-zero Fourier
  transform, translates of $\varphi$ span a dense subspace of $L^1$.
\end{theorem}

\begin{theorem}[Rudin 9.5, closed translation-invariant ideals]
  \label{thm:closed-invariant-ideals-L1}
  \uses{thm:wiener-lemma}
  \lean{Rudin.Pending.closed_invariant_ideals_L1}
  \leanok
  Closed translation-invariant ideals in $L^1(\mathbb R^n)$ are in
  bijection with closed subsets of $\mathbb R^n$ via the spectrum
  (hull of the ideal).
\end{theorem}

\begin{theorem}[Rudin 9.7, Wiener's Tauberian theorem]
  \label{thm:wiener-tauberian}
  \uses{thm:closed-invariant-ideals-L1}
  \lean{Rudin.Pending.wiener_tauberian}
  \leanok
  Let $K \in L^1(\mathbb R^n)$ have nowhere-zero Fourier transform.
  If $f \in L^\infty$ satisfies $(K \star f)(x) \to A \int K$, then
  $(h \star f)(x) \to A \int h$ for every $h \in L^1$.
\end{theorem}

\begin{theorem}[Rudin 9.10, Ikehara]
  \label{thm:ikehara}
  \uses{thm:wiener-tauberian}
  \lean{Rudin.Pending.ikehara}
  \leanok
  Asymptotic distribution theorem via analytic continuation of
  Dirichlet series, classically used to derive the prime number
  theorem.
\end{theorem}

\begin{theorem}[Rudin 9.12, Prime Number Theorem]
  \label{thm:pnt}
  \uses{thm:ikehara}
  \lean{Rudin.Pending.prime_number_theorem}
  \leanok
  $\pi(x) \sim x / \log x$.
\end{theorem}

\chapter{Banach Algebras}

\begin{theorem}[Rudin 10.11, units form an open set]
  \label{thm:units-open}
  \lean{Rudin.Ch10.units_isOpen}
  \leanok
  In a Banach algebra, the set of invertible elements is open.
\end{theorem}
\begin{proof}
  \leanok
\end{proof}

\begin{theorem}[Rudin 10.12, inversion is continuous]
  \label{thm:inversion-continuous}
  \uses{thm:units-open}
  \lean{Rudin.Ch10.inverse_continuousAt}
  \leanok
  The inversion map $x \mapsto x^{-1}$ is continuous at every unit.
\end{theorem}
\begin{proof}
  \leanok
\end{proof}

\begin{definition}[Rudin 10.7, Neumann unit construction]
  \label{def:neumann-unit}
  \lean{Rudin.Ch10.oneSubUnit}
  \leanok
  For $\|a\| < 1$, $1 - a$ is explicitly a unit, witnessed by the
  geometric series $\sum_{n \ge 0} a^n$.
\end{definition}

\begin{theorem}[Rudin 10.7, Neumann series formula]
  \label{thm:neumann-series}
  \uses{def:neumann-unit}
  \lean{Rudin.Ch10.neumann_series}
  \leanok
  For $\|a\| < 1$ in a Banach algebra,
  $\sum_{n \ge 0} a^n = (1 - a)^{-1}$.
\end{theorem}
\begin{proof}
  \leanok
\end{proof}

\begin{theorem}[$1 - a$ is a unit for small $a$]
  \label{thm:one-sub-unit}
  \uses{thm:neumann-series}
  \lean{Rudin.Ch10.isUnit_one_sub}
  \leanok
  If $\|a\| < 1$ in a Banach algebra, then $1 - a$ is invertible.
\end{theorem}
\begin{proof}
  \leanok
\end{proof}

\begin{theorem}[Rudin 10.13, spectrum nonempty]
  \label{thm:spectrum-nonempty}
  \lean{Rudin.Ch10.spectrum_nonempty}
  \leanok
  Every element of a nontrivial complex Banach algebra has nonempty spectrum.
\end{theorem}
\begin{proof}
  \leanok
\end{proof}

\begin{theorem}[Rudin 10.14, Gelfand--Mazur]
  \label{thm:gelfand-mazur}
  \uses{thm:spectrum-nonempty}
  \lean{Rudin.Ch10.gelfandMazur}
  \leanok
  If every nonzero element of a complex Banach algebra $A$ is invertible,
  then the natural map $\mathbb{C} \to A$ is an isometric algebra
  isomorphism.
\end{theorem}
\begin{proof}
  \leanok
\end{proof}

\begin{theorem}[Rudin 10.13, Gelfand's spectral radius formula]
  \label{thm:gelfand-formula}
  \lean{Rudin.Ch10.gelfand_formula}
  \leanok
  For $a$ in a complex Banach algebra,
  $\|a^n\|^{1/n} \to \rho(a)$ as $n \to \infty$,
  where $\rho(a)$ is the spectral radius.
\end{theorem}
\begin{proof}
  \leanok
\end{proof}

\begin{theorem}[Rudin 10.13, spectrum bounded by norm]
  \label{thm:spectrum-in-closed-ball}
  \lean{Rudin.Ch10.spectrum_subset_closedBall}
  \leanok
  $\sigma(a) \subseteq \{z : |z| \le \|a\|\}$.
\end{theorem}
\begin{proof}
  \leanok
\end{proof}

\begin{theorem}[Rudin 10.13, spectrum is compact]
  \label{thm:spectrum-compact}
  \uses{thm:spectrum-in-closed-ball}
  \lean{Rudin.Ch10.spectrum_isCompact}
  \leanok
  The spectrum of any element of a complex Banach algebra is compact.
\end{theorem}
\begin{proof}
  \leanok
\end{proof}

\begin{theorem}[resolvent set is open]
  \label{thm:resolventSet-open}
  \lean{Rudin.Ch10.resolventSet_isOpen}
  \leanok
  The complement of the spectrum in $\mathbb{C}$ is open.
\end{theorem}
\begin{proof}
  \leanok
\end{proof}

\begin{theorem}[spectral radius $\le$ norm]
  \label{thm:spectralRadius-le-norm}
  \uses{thm:spectrum-in-closed-ball}
  \lean{Rudin.Ch10.spectralRadius_le_norm}
  \leanok
  $\rho(a) \le \|a\|$ for any element of a complex Banach algebra.
\end{theorem}
\begin{proof}
  \leanok
\end{proof}

\begin{theorem}[corollary: $\rho(1) \le 1$]
  \label{thm:spectralRadius-one-le}
  \uses{thm:spectralRadius-le-norm}
  \lean{Rudin.Ch10.spectralRadius_one_le}
  \leanok
  The spectral radius of the identity element is at most $1$.
  \emph{Proved in-project from $\rho(a) \le \|a\|$ and $\|1\| = 1$.}
\end{theorem}
\begin{proof}
  \leanok
\end{proof}

% --- Ch10 Pending ---

\begin{theorem}[Rudin 10.10, group of units]
  \label{thm:units-def}
  \lean{Rudin.Pending.units_def}
  \leanok
  Definition of $G(A)$, the group of invertible elements; mathlib
  packages this as $A^\times$.
\end{theorem}

\begin{theorem}[Rudin 10.18, spectrum upper semicontinuity]
  \label{thm:spectrum-uppersemiContinuous}
  \lean{Rudin.Pending.spectrum_upperSemiContinuous}
  \leanok
  If $x_n \to x$ then $\sigma(x_n) \to \sigma(x)$ in a suitable
  Hausdorff sense.
\end{theorem}

\begin{theorem}[Rudin 10.19, polynomial spectral mapping]
  \label{thm:spectrum-polynomial}
  \lean{Rudin.Pending.spectrum_polynomial_mapping}
  \leanok
  $\sigma(p(x)) = p(\sigma(x))$ for any polynomial $p$.
\end{theorem}

\begin{theorem}[Rudin 10.20, resolvent is holomorphic]
  \label{thm:resolvent-holomorphic}
  \lean{Rudin.Pending.resolvent_holomorphic}
  \leanok
  The resolvent function $R(\lambda) = (\lambda - x)^{-1}$ is
  holomorphic on the resolvent set.
\end{theorem}

\begin{theorem}[Rudin 10.21--10.27, holomorphic functional calculus]
  \label{thm:holomorphic-fc}
  \uses{thm:resolvent-holomorphic}
  \lean{Rudin.Pending.holomorphic_functional_calculus}
  \leanok
  If $f$ is holomorphic on a neighbourhood of $\sigma(x)$, then
  $f(x)$ is well defined via Cauchy's integral formula.
\end{theorem}

\begin{theorem}[Rudin 10.28, holomorphic spectral mapping]
  \label{thm:spectrum-holomorphic}
  \uses{thm:holomorphic-fc}
  \lean{Rudin.Pending.spectrum_holomorphic_mapping}
  \leanok
  $\sigma(f(x)) = f(\sigma(x))$ for $f$ holomorphic.
\end{theorem}

\begin{theorem}[Rudin 10.30, Shilov idempotent]
  \label{thm:shilov-idempotent}
  \uses{thm:holomorphic-fc}
  \lean{Rudin.Pending.shilov_idempotent}
  \leanok
  If $\sigma(x) = K_1 \sqcup K_2$ with $K_i$ disjoint compact, there
  exist orthogonal idempotents $e_i \in A$ with $e_1 + e_2 = 1$ and
  $\sigma(e_i) \subseteq K_i \cup \{0\}$.
\end{theorem}

\begin{theorem}[Rudin 10.33, square root]
  \label{thm:banachAlgebra-sqrt}
  \uses{thm:holomorphic-fc}
  \lean{Rudin.Pending.banachAlgebra_squareRoot}
  \leanok
  Every $x$ with $\sigma(x) \cap (-\infty, 0] = \emptyset$ admits a
  holomorphic square root.
\end{theorem}

\begin{theorem}[Rudin 10.34, logarithm]
  \label{thm:banachAlgebra-log}
  \uses{thm:banachAlgebra-sqrt}
  \lean{Rudin.Pending.banachAlgebra_log}
  \leanok
  Every $x$ with $\sigma(x) \cap (-\infty, 0] = \emptyset$ admits a
  logarithm.
\end{theorem}

\chapter{Commutative Banach Algebras}

\begin{definition}[Rudin 11.5, character space]
  \label{def:character-space}
  \lean{Rudin.Ch11.CharacterSpace}
  \leanok
  The character space $\Delta$ of a commutative complex Banach algebra
  $A$ is the set of continuous nonzero algebra homomorphisms $A \to \mathbb{C}$,
  equipped with the weak-$*$ topology.
\end{definition}

\begin{theorem}[Rudin 11.9, characters detect the spectrum]
  \label{thm:spectrum-via-characters}
  \uses{def:character-space, thm:spectrum-nonempty}
  \lean{Rudin.Ch11.mem_spectrum_iff_exists_character}
  \leanok
  $z \in \sigma(a)$ iff there is a character $\varphi \in \Delta$ with
  $\varphi(a) = z$.
\end{theorem}
\begin{proof}
  \leanok
\end{proof}

\begin{definition}[Rudin 11.18, Gelfand transform]
  \label{def:gelfand-transform}
  \uses{def:character-space}
  \lean{Rudin.Ch11.gelfand}
  \leanok
  The Gelfand transform $A \to C(\Delta, \mathbb{C})$ sends $a$ to
  $\varphi \mapsto \varphi(a)$.
\end{definition}

\begin{theorem}[Rudin 11.18, Gelfand--Naimark: commutative C$^{*}$-algebras]
  \label{thm:gelfand-naimark}
  \uses{def:gelfand-transform, thm:spectrum-via-characters}
  \lean{Rudin.Ch11.gelfand_isometry, Rudin.Ch11.gelfand_bijective}
  \leanok
  For a commutative unital C$^{*}$-algebra $A$, the Gelfand transform
  is an isometric $*$-isomorphism $A \cong C(\Delta, \mathbb{C})$.
\end{theorem}
\begin{proof}
  \leanok
\end{proof}

\begin{definition}[Rudin 11.5, character from a maximal ideal]
  \label{def:maximal-ideal-character}
  \uses{def:character-space, thm:gelfand-mazur}
  \lean{Rudin.Ch11.maximalIdealCharacter}
  \leanok
  Every maximal ideal $I \subseteq A$ determines a character via the
  composition $A \twoheadrightarrow A/I \xrightarrow{\cong} \mathbb{C}$.
\end{definition}

\begin{theorem}[character vanishes on its ideal]
  \label{thm:max-ideal-char-zero}
  \uses{def:maximal-ideal-character}
  \lean{Rudin.Ch11.maximalIdealCharacter_apply_eq_zero_of_mem}
  \leanok
  The character of a maximal ideal $I$ vanishes on every $a \in I$.
\end{theorem}
\begin{proof}
  \leanok
\end{proof}

% --- Ch11 Pending ---

\begin{theorem}[Rudin 11.7, character / maximal-ideal correspondence]
  \label{thm:character-maxIdeal-bijection}
  \uses{def:character-space, def:maximal-ideal-character}
  \lean{Rudin.Pending.character_maxIdeal_bijection}
  \leanok
  The character space $\Delta(A)$ is in bijection with the maximal
  ideals of a commutative unital Banach algebra $A$.
\end{theorem}

\begin{theorem}[Rudin 11.13, $\|\widehat{x}\|_\infty = \rho(x)$]
  \label{thm:gelfand-norm-spectralRadius}
  \uses{def:gelfand-transform, thm:gelfand-formula}
  \lean{Rudin.Pending.gelfand_norm_eq_spectralRadius}
  \leanok
  $\|\widehat{x}\|_\infty = \rho(x)$ for the Gelfand transform.
\end{theorem}

\begin{theorem}[Rudin 11.18, continuous functional calculus]
  \label{thm:cfc-commutative}
  \uses{thm:gelfand-naimark}
  \lean{Rudin.Pending.continuous_functional_calculus_commutative}
  \leanok
  Continuous functional calculus on commutative unital
  $C^*$-algebras.
\end{theorem}

\begin{theorem}[Rudin 11.21, Stone--Čech via characters]
  \label{thm:stoneCech-via-characters}
  \uses{def:character-space}
  \lean{Rudin.Pending.stoneCech_via_characters}
  \leanok
  For a Tychonoff space $X$, the Stone--Čech compactification $\beta
  X$ arises as the character space of $C_b(X)$.
\end{theorem}

\begin{theorem}[Rudin 11.23, existence of involution]
  \label{thm:commutative-involution}
  \lean{Rudin.Pending.commutative_involution_existence}
  \leanok
  Existence of an involution on commutative semisimple Banach
  algebras under suitable conditions.
\end{theorem}

\chapter{Bounded Operators on a Hilbert Space}

\begin{definition}[Rudin 12.9, adjoint of a bounded operator]
  \label{def:adjoint}
  \lean{Rudin.Ch12.adjoint}
  \leanok
  For Hilbert spaces $E$, $F$ and $A \in \mathcal{B}(E,F)$, there is a
  unique bounded operator $A^{*} \in \mathcal{B}(F,E)$ characterised
  by $\langle A x, y\rangle = \langle x, A^{*} y\rangle$ for all
  $x \in E$, $y \in F$.
\end{definition}

\begin{theorem}[Rudin 12.9, defining relation of the adjoint]
  \label{thm:adjoint-inner-left}
  \uses{def:adjoint}
  \lean{Rudin.Ch12.adjoint_inner_left}
  \leanok
  $\langle A^{*} y, x\rangle = \langle y, A x\rangle$.
\end{theorem}
\begin{proof}
  \leanok
\end{proof}

\begin{theorem}[Rudin 12.9, adjoint is involutive]
  \label{thm:adjoint-adjoint}
  \uses{def:adjoint}
  \lean{Rudin.Ch12.adjoint_adjoint}
  \leanok
  $(A^{*})^{*} = A$.
\end{theorem}
\begin{proof}
  \leanok
\end{proof}

\begin{theorem}[Rudin 12.10, adjoint of composition]
  \label{thm:adjoint-comp}
  \uses{def:adjoint}
  \lean{Rudin.Ch12.adjoint_comp}
  \leanok
  $(AB)^{*} = B^{*} A^{*}$.
\end{theorem}
\begin{proof}
  \leanok
\end{proof}

\begin{theorem}[Rudin 12.11, C$^{*}$-identity]
  \label{thm:cstar-identity}
  \uses{def:adjoint}
  \lean{Rudin.Ch12.norm_adjoint_comp_self}
  \leanok
  $\|A^{*} A\| = \|A\|^{2}$.
\end{theorem}
\begin{proof}
  \leanok
\end{proof}

\begin{theorem}[in-project corollaries of 12.9]
  \label{thm:adjoint-linear-laws}
  \uses{def:adjoint}
  \lean{Rudin.Ch12.adjoint_zero, Rudin.Ch12.adjoint_add, Rudin.Ch12.adjoint_neg, Rudin.Ch12.adjoint_sub, Rudin.Ch12.adjoint_norm}
  \leanok
  $0^{*} = 0$, $(A + B)^{*} = A^{*} + B^{*}$,
  $(-A)^{*} = -A^{*}$, $(A - B)^{*} = A^{*} - B^{*}$, and
  $\|A^{*}\| = \|A\|$.
  \emph{Proved in-project from the defining inner-product relation.}
\end{theorem}
\begin{proof}
  \leanok
\end{proof}

\begin{theorem}[in-project: $A^{*} A$ is self-adjoint]
  \label{thm:adjoint-self-comp}
  \uses{thm:adjoint-comp, thm:self-adjoint-iff}
  \lean{Rudin.Ch12.isSelfAdjoint_adjoint_comp_self}
  \leanok
  For any bounded operator $A\colon E \to F$, $A^{*}A\colon E \to E$ is
  self-adjoint. \emph{Proved in-project by
  $(A^{*}A)^{*} = A^{*}(A^{*})^{*} = A^{*}A$.}
\end{theorem}
\begin{proof}
  \leanok
\end{proof}

\begin{theorem}[Rudin 12.12, self-adjoint iff $A^{*} = A$]
  \label{thm:self-adjoint-iff}
  \uses{def:adjoint}
  \lean{Rudin.Ch12.isSelfAdjoint_iff_adjoint_eq}
  \leanok
  $A$ is self-adjoint iff $A^{*} = A$.
\end{theorem}
\begin{proof}
  \leanok
\end{proof}

\begin{theorem}[Rudin 12.23, eigenvalues of a self-adjoint operator are real]
  \label{thm:self-adjoint-real-eigenvalue}
  \uses{thm:self-adjoint-iff}
  \lean{Rudin.Ch12.self_adjoint_eigenvalue_real}
  \leanok
  Every eigenvalue $\mu$ of a self-adjoint operator satisfies
  $\overline{\mu} = \mu$.
\end{theorem}
\begin{proof}
  \leanok
\end{proof}

\begin{theorem}[Rudin 12.26, characterisation of normal operators]
  \label{thm:normal-iff}
  \uses{def:adjoint}
  \lean{Rudin.Ch12.isStarNormal_iff_norm_eq_adjoint}
  \leanok
  An operator $T$ on a Hilbert space is normal ($T^{*}T = T T^{*}$)
  iff $\|T v\| = \|T^{*} v\|$ for every $v$.
\end{theorem}
\begin{proof}
  \leanok
\end{proof}

\begin{theorem}[Rudin 12.26, kernel $=$ range orthogonal complement for normal operators]
  \label{thm:normal-orthogonal-range}
  \uses{thm:normal-iff}
  \lean{Rudin.Ch12.normal_orthogonal_range}
  \leanok
  For a normal operator $T$, $(\operatorname{range} T)^{\perp} = \ker T$.
\end{theorem}
\begin{proof}
  \leanok
\end{proof}

\begin{definition}[orthogonal projection onto a closed subspace]
  \label{def:ortho-proj}
  \lean{Rudin.Ch12.orthogonalProjection}
  \leanok
  For a closed subspace $U \subseteq E$ admitting an orthogonal
  projection, $P_U\colon E \to E$ projects onto $U$ along $U^{\perp}$.
\end{definition}

\begin{theorem}[Rudin 12.14, orthogonal projection is self-adjoint]
  \label{thm:ortho-proj-self-adjoint}
  \uses{def:ortho-proj, thm:self-adjoint-iff}
  \lean{Rudin.Ch12.orthogonalProjection_isSelfAdjoint}
  \leanok
  Every orthogonal projection $P_U$ satisfies $P_U^{*} = P_U$.
\end{theorem}
\begin{proof}
  \leanok
\end{proof}

\begin{theorem}[Rudin 12.22, spectral theorem (finite-dimensional)]
  \label{thm:spectral-diagonalization}
  \uses{def:adjoint}
  \lean{Rudin.Ch12.spectralDiagonalization, Rudin.Ch12.spectralDiagonalization_apply_self_apply}
  \leanok
  A self-adjoint operator $T$ on a finite-dimensional inner product
  space $E$ is unitarily equivalent to the diagonal operator on the
  orthogonal direct sum of its eigenspaces: there is a linear isometry
  $E \cong \bigoplus_\mu E_\mu$ conjugating $T$ to multiplication by
  $\mu$ on each summand.
\end{theorem}
\begin{proof}
  \leanok
\end{proof}

% --- Ch12 Pending ---

\begin{theorem}[Rudin 12.13, polarisation identity]
  \label{thm:polarisation}
  \lean{Rudin.Pending.polarisation}
  \leanok
  The inner product is recovered from the norm via polarisation.
\end{theorem}

\begin{theorem}[Rudin 12.16, polar decomposition]
  \label{thm:polar-decomposition}
  \uses{def:adjoint}
  \lean{Rudin.Pending.polar_decomposition}
  \leanok
  $T = U |T|$ with $U$ partial isometry and $|T| = (T^* T)^{1/2}$.
\end{theorem}

\begin{theorem}[Rudin 12.17, isometry characterisation]
  \label{thm:isometry-iff}
  \uses{def:adjoint}
  \lean{Rudin.Pending.isometry_iff_adjoint_id}
  \leanok
  $T$ is an isometry iff $T^* T = I$.
\end{theorem}

\begin{theorem}[Rudin 12.18, self-adjoint via real quadratic form]
  \label{thm:selfAdjoint-realQuadratic}
  \uses{thm:self-adjoint-iff}
  \lean{Rudin.Pending.selfAdjoint_iff_realQuadratic}
  \leanok
  $T = T^*$ iff $\langle T x, x \rangle \in \mathbb R$ for all $x$
  (over $\mathbb C$).
\end{theorem}

\begin{theorem}[Rudin 12.19--12.20, continuous functional calculus]
  \label{thm:cfc-selfAdjoint}
  \uses{thm:self-adjoint-iff}
  \lean{Rudin.Pending.cfc_selfAdjoint}
  \leanok
  Continuous functional calculus for self-adjoint operators extends
  polynomial calculus to $C(\sigma(T))$.
\end{theorem}

\begin{theorem}[Rudin 12.21, multiplication-operator spectral theorem]
  \label{thm:spectral-multiplication}
  \uses{thm:cfc-selfAdjoint}
  \lean{Rudin.Pending.spectral_theorem_normal}
  \leanok
  Every normal operator is unitarily equivalent to multiplication by
  a measurable function on some $L^2(\mu)$.
\end{theorem}

\begin{theorem}[Rudin 12.24, projection-valued spectral resolution]
  \label{thm:spectral-resolution}
  \uses{thm:spectral-multiplication}
  \lean{Rudin.Pending.spectral_resolution}
  \leanok
  Every normal $T$ admits a unique resolution of the identity $E$ on
  $\sigma(T)$ with $T = \int \lambda \, dE(\lambda)$.
\end{theorem}

\begin{theorem}[Rudin 12.25, Borel functional calculus]
  \label{thm:borel-fc}
  \uses{thm:spectral-resolution}
  \lean{Rudin.Pending.borel_functional_calculus}
  \leanok
  $f(T) = \int f \, dE$ for bounded Borel $f$ on $\sigma(T)$.
\end{theorem}

\begin{theorem}[Rudin 12.27, positive operators]
  \label{thm:positive-iff-quadratic}
  \uses{def:adjoint}
  \lean{Rudin.Pending.positive_iff_quadratic_nonneg}
  \leanok
  $T \ge 0$ iff $\langle T x, x \rangle \ge 0$, iff $T = S^* S$.
\end{theorem}

\begin{theorem}[Rudin 12.30, spectrum of normal $=$ essential range]
  \label{thm:spectrum-normal-essRange}
  \uses{thm:spectral-multiplication}
  \lean{Rudin.Pending.spectrum_normal_essRange}
  \leanok
  Spectrum of a normal operator equals the essential range of its
  symbol under the multiplication-operator representation.
\end{theorem}

\begin{theorem}[Rudin 12.32, commutant of normal operator]
  \label{thm:commutant-normal}
  \uses{thm:spectral-resolution}
  \lean{Rudin.Pending.commutant_normal}
  \leanok
  The commutant of a normal operator equals the commutant of its
  spectral measure.
\end{theorem}

\begin{theorem}[Rudin 12.36, multiplicity theory]
  \label{thm:multiplicity-normal}
  \uses{thm:spectral-resolution}
  \lean{Rudin.Pending.multiplicity_normal}
  \leanok
  Multiplicity theory for normal operators.
\end{theorem}

\begin{theorem}[Rudin 12.40, mean ergodic theorem]
  \label{thm:mean-ergodic}
  \lean{Rudin.Pending.mean_ergodic_vonNeumann}
  \leanok
  For an isometry $T$ on a Hilbert space, $\frac{1}{N} \sum_{k=0}^{N-1}
  T^k x \to P x$ strongly, with $P$ projecting on $\ker(T - I)$.
\end{theorem}

\begin{theorem}[Rudin 12.43, individual ergodic theorem (Hilbert)]
  \label{thm:pointwise-ergodic-hilbert}
  \uses{thm:mean-ergodic}
  \lean{Rudin.Pending.pointwise_ergodic_hilbert}
  \leanok
  Pointwise / individual ergodic theorem in Hilbert-space form.
\end{theorem}

\chapter{Unbounded Operators}

\begin{definition}[Rudin 13.1, closed operator]
  \label{def:closed-operator}
  \lean{Rudin.Ch13.IsClosedOperator}
  \leanok
  A densely-defined linear operator $T$ is \emph{closed} iff its graph
  $\{(x, T x) : x \in \operatorname{dom} T\}$ is closed in $E \times F$.
\end{definition}

\begin{definition}[Rudin 13.4, closable operator]
  \label{def:closable-operator}
  \uses{def:closed-operator}
  \lean{Rudin.Ch13.IsClosableOperator}
  \leanok
  An operator $T$ is \emph{closable} iff the closure of its graph is
  itself the graph of some (uniquely determined) operator.
\end{definition}

\begin{theorem}[Closed operators are closable]
  \label{thm:closed-closable}
  \uses{def:closed-operator, def:closable-operator}
  \lean{Rudin.Ch13.isClosable_of_isClosed}
  \leanok
  Every closed operator is trivially closable.
\end{theorem}
\begin{proof}
  \leanok
\end{proof}

\begin{definition}[Rudin 13.6, closure of an operator]
  \label{def:closure-operator}
  \uses{def:closable-operator}
  \lean{Rudin.Ch13.closure}
  \leanok
  If $T$ is closable, $\overline{T}$ is the unique extension whose
  graph is the closure of the graph of $T$.
\end{definition}

\begin{theorem}[Rudin 13.7, the closure is closed]
  \label{thm:closure-closed}
  \uses{def:closure-operator}
  \lean{Rudin.Ch13.closure_isClosed}
  \leanok
  If $T$ is closable, then $\overline{T}$ is a closed operator.
\end{theorem}
\begin{proof}
  \leanok
\end{proof}

\begin{definition}[Rudin 13.9, adjoint of an unbounded operator]
  \label{def:pmap-adjoint}
  \uses{def:closed-operator}
  \lean{Rudin.Ch13.pmapAdjoint}
  \leanok
  For a densely-defined operator $T\colon E \to F$ between Hilbert
  spaces, its adjoint $T^{*}$ is defined on the subspace of $y \in F$
  for which $x \mapsto \langle T x, y\rangle$ is bounded on
  $\operatorname{dom} T$.
\end{definition}

\begin{theorem}[Rudin 13.10, the adjoint is closed]
  \label{thm:pmap-adjoint-closed}
  \uses{def:pmap-adjoint, def:closed-operator}
  \lean{Rudin.Ch13.pmapAdjoint_isClosed}
  \leanok
  The adjoint of any densely-defined operator is a closed operator.
\end{theorem}
\begin{proof}
  \leanok
\end{proof}

\begin{theorem}[Rudin 13.11, adjoint is a formal adjoint]
  \label{thm:pmap-adjoint-formal}
  \uses{def:pmap-adjoint}
  \lean{Rudin.Ch13.pmapAdjoint_isFormalAdjoint}
  \leanok
  For a densely-defined $T$, $\langle T x, y\rangle = \langle x,
  T^{*} y\rangle$ whenever $x \in \operatorname{dom} T$ and $y \in
  \operatorname{dom} T^{*}$.
\end{theorem}
\begin{proof}
  \leanok
\end{proof}

% --- Ch13 Pending ---

\begin{theorem}[Rudin 13.13, adjoint dense iff closable]
  \label{thm:adjoint-dense-iff-closable}
  \uses{def:closable-operator, def:pmap-adjoint}
  \lean{Rudin.Pending.adjoint_denselyDefined_iff_closable}
  \leanok
  The adjoint of a closed densely-defined operator is densely defined
  iff the operator is closable.
\end{theorem}

\begin{theorem}[Rudin 13.14, $T^{**} = \overline{T}$]
  \label{thm:adjoint-adjoint-closure}
  \uses{def:closure-operator, def:pmap-adjoint}
  \lean{Rudin.Pending.adjoint_adjoint_eq_closure}
  \leanok
  The double adjoint of a densely-defined closable operator equals
  its closure.
\end{theorem}

\begin{theorem}[Rudin 13.16, Cayley transform]
  \label{thm:cayley-transform}
  \lean{Rudin.Pending.cayley_transform}
  \leanok
  Closed symmetric $T$ corresponds to an isometry
  $U = (T - i)(T + i)^{-1}$.
\end{theorem}

\begin{theorem}[Rudin 13.17, self-adjoint iff Cayley unitary]
  \label{thm:selfAdjoint-cayley-unitary}
  \uses{thm:cayley-transform}
  \lean{Rudin.Pending.selfAdjoint_iff_cayley_unitary}
  \leanok
  $T$ is self-adjoint iff its Cayley transform is unitary.
\end{theorem}

\begin{theorem}[Rudin 13.19, deficiency indices]
  \label{thm:deficiency-indices}
  \uses{thm:cayley-transform}
  \lean{Rudin.Pending.selfAdjoint_extension_iff_deficiency}
  \leanok
  Self-adjoint extensions of a closed symmetric $T$ exist iff
  $\dim \ker(T^* - i) = \dim \ker(T^* + i)$.
\end{theorem}

\begin{theorem}[Rudin 13.22, spectral theorem (unbounded self-adjoint)]
  \label{thm:spectral-unbounded-selfAdjoint}
  \uses{thm:selfAdjoint-cayley-unitary, thm:spectral-resolution}
  \lean{Rudin.Pending.spectral_theorem_unbounded_selfAdjoint}
  \leanok
  Every self-adjoint operator $T$ admits a unique spectral resolution
  $T = \int \lambda \, dE(\lambda)$ over $\mathbb R$.
\end{theorem}

\begin{theorem}[Rudin 13.24, unbounded functional calculus]
  \label{thm:unbounded-fc}
  \uses{thm:spectral-unbounded-selfAdjoint}
  \lean{Rudin.Pending.unbounded_functional_calculus}
  \leanok
  Functional calculus for unbounded self-adjoint operators, defined
  via the spectral resolution and including unbounded measurable
  functions.
\end{theorem}

\begin{theorem}[Rudin 13.30, spectral theorem (unbounded normal)]
  \label{thm:spectral-unbounded-normal}
  \uses{thm:spectral-unbounded-selfAdjoint}
  \lean{Rudin.Pending.spectral_theorem_unbounded_normal}
  \leanok
  Spectral theorem for unbounded normal operators.
\end{theorem}

\begin{theorem}[Rudin 13.33, Stone's theorem]
  \label{thm:stone-theorem}
  \uses{thm:spectral-unbounded-selfAdjoint}
  \lean{Rudin.Pending.stone_theorem}
  \leanok
  A strongly continuous one-parameter unitary group has the form
  $U(t) = e^{itA}$ for a unique self-adjoint $A$.
\end{theorem}

\begin{theorem}[Rudin 13.36, self-adjoint contraction generators]
  \label{thm:selfAdjoint-contraction}
  \uses{thm:stone-theorem}
  \lean{Rudin.Pending.selfAdjoint_generator_contraction}
  \leanok
  Self-adjoint generators of contraction semigroups (Hille--Yosida
  special case).
\end{theorem}

\begin{theorem}[Rudin 13.37--13.38, Friedrichs extension]
  \label{thm:friedrichs}
  \lean{Rudin.Pending.friedrichs_extension}
  \leanok
  A densely-defined semibounded symmetric operator has a canonical
  self-adjoint extension.
\end{theorem}
